\documentclass[11pt]{ctexart}
\usepackage[margin=2.5cm]{geometry}
\usepackage{amsmath,amssymb}
\usepackage{booktabs,tabularx}
\usepackage[authoryear,round]{natbib}
\usepackage[colorlinks=true,linkcolor=blue,citecolor=blue,urlcolor=blue]{hyperref}

\newcommand{\E}{\mathbb{E}}
\newcommand{\F}{\mathcal{F}}
\newcommand{\dist}{\operatorname{dist}}

\begin{document}

\section{JOTA 讲了什么}

\paragraph{Summary.}
\begin{itemize}
  \item \textbf{机制：}一条能量不等式，噪声通过 $s_k^2t_k^2\sigma_k^2$ 进入；它与 SGD 的偏差--方差结构同构。
  \item \textbf{速率：}噪声递减时恢复 a.s.\ $o(1/k^2)$；持续噪声时最好约为 $1/\sqrt{k}$，与 SGD 同阶。
  \item \textbf{迭代点：}保住加速的体制下得到 a.s.\ 收敛，但需要 $\ell^1$ 型噪声条件和不降的步长。
  % \item \textbf{未回答的问题：}这些条件是否紧；持续噪声下的迭代点结论；临界 $\alpha=3$。
\end{itemize}

\citet{attouch2026ravine} 用一条随机 Lyapunov 不等式，统一分析一般外推参数 $\alpha_k$ 下的随机 Nesterov（SNAG）与 Ravine（SRAG）方法。噪声衰减得足够快时，恢复确定性的 $o(1/k^2)$ 与迭代点收敛；噪声持续时，最好速率约为 $1/\sqrt{k}$，与 SGD 同阶。

\begin{table}[htbp]
    \centering
    \small
    \renewcommand{\arraystretch}{1.25}
    \caption{随机动量方法与 JOTA 论文贡献的对比}
    \label{tab:jota-contribution}
    \begin{tabularx}{\textwidth}{
        @{}
        >{\raggedright\arraybackslash}p{2.2cm}
        >{\raggedright\arraybackslash}p{3.0cm}
        >{\raggedright\arraybackslash}p{3.0cm}
        >{\raggedright\arraybackslash}X
        @{}
    }
        \toprule
        比较维度
        & Sebbouh et al.~(2021)
        & Liu--Yuan (2022)
        & JOTA (2026) \\
        \midrule

        算法
        & SHB
        & SHB、SNAG
        & SNAG、SRAG 统一分析 \\

        动量参数
        & $\beta_k$ 由步长决定
        & 常数 $\bar{\alpha}\in(0,1)$
        & 一般 $\alpha_k$，包括 $\alpha_k\to1$ \\

        动量阶数
        & 与步长绑定
        & $O(1)$
        & $\alpha_k=1-\alpha/k$，
          $1-\alpha_k=\Theta(k^{-1})$，
          $t_k=\Theta(k)$ \\

        步长
        & $\eta_k\asymp k^{-1/2-\delta}$
        & $s_k\asymp k^{-2/3-\delta}$
        & 可取常数，也可递减 \\

        噪声体制
        & 持续噪声
        & 持续噪声
        & 噪声递减 \\

        函数值速率
        & 接近 $o(k^{-1/2})$
        & $O(k^{-1/3+\delta})$ a.s.
        & $\alpha>3$ 时恢复 $o(1/k^2)$ \\

        迭代点收敛
        & 有限维 a.s. 收敛
        & 有限维 a.s. 收敛
        & 有限维强 a.s. 收敛，
          无限维弱 a.s. 收敛 \\

        Ravine 方法
        & --
        & --
        & 通过 $\gamma_k=\alpha_{k+1}$
          转移 SNAG 结论 \\

        主要贡献
        & 持续噪声下的 SHB 理论
        & 常数动量下的 SNAG/SHB 理论
        & 噪声递减、一般 $\alpha_k$、
          加速速率与 Ravine 统一分析 \\

        \bottomrule
    \end{tabularx}
\end{table}

\subsection{设定}

\paragraph{算法与噪声.}
\[
y_k=x_k+\alpha_k(x_k-x_{k-1}),
\qquad
x_{k+1}=y_k-s_k\bigl(\nabla f(y_k)+e_k\bigr),
\]
\[
\E[e_k\mid\F_k]=0,
\qquad
\E\bigl[\|e_k\|^2\mid\F_k\bigr]\le\sigma_k^2 .
\]

\paragraph{$t_k$ 统一所有动量.}
令
\[
t_k:=1+\sum_{i\ge k}\prod_{j=k}^{i}\alpha_j ,
\]
其中 (K0) 保证级数有限，则
\[
\alpha_k=\frac{t_k-1}{t_{k+1}} .
\]
对 $\alpha_k=1-\alpha/k$，有 $t_k\asymp k/(\alpha-1)$。这一参数化来自确定性惯性算法的分析 \citep{attouch2018cabot}。

\paragraph{分析变量 $z_k$.}
令
\[
z_k:=x_{k-1}+t_k(x_k-x_{k-1}),
\]
则
\[
z_{k+1}=z_k-s_kt_{k+1}\bigl(\nabla f(y_k)+e_k\bigr),
\]
\[
x_k=\Bigl(1-\frac1{t_k}\Bigr)x_{k-1}+\frac1{t_k}z_k,
\qquad
y_k=\Bigl(1-\frac1{t_{k+1}}\Bigr)x_k+\frac1{t_{k+1}}z_k .
\]
\textbf{$z_k$ 是步长为 $s_kt_{k+1}$ 的类 SGD 序列，$x_k$ 是它的在线平均，$y_k$ 是两者的凸组合。}

$z_k$ 的递推一行即可验证。由定义 $z_k=x_k+(t_k-1)(x_k-x_{k-1})$，且 $t_{k+1}\alpha_k=t_k-1$。代入
\[
z_{k+1}=x_k+t_{k+1}(x_{k+1}-x_k)
=x_k+t_{k+1}(y_k-x_k)-s_kt_{k+1}\bigl(\nabla f(y_k)+e_k\bigr),
\]
而 $t_{k+1}(y_k-x_k)=t_{k+1}\alpha_k(x_k-x_{k-1})=(t_k-1)(x_k-x_{k-1})$，所以
\[
z_{k+1}=z_k-s_kt_{k+1}\bigl(\nabla f(y_k)+e_k\bigr).
\]

\paragraph{动量条件.}
(K1)
\[
t_{k+1}^2-t_k^2\le t_{k+1}
\]
对应 $\alpha\ge3$；(K1$^+$) 把右端换成 $m\,t_{k+1}$、$m<1$，对应 $\alpha>3$。在连续时间中，$\alpha\ge3$ 给出 $O(1/t^2)$ \citep{su2016ode}，$\alpha>3$ 给出 $o(1/t^2)$ 与轨道收敛 \citep{attouch2018fast}；离散的 $o(1/k^2)$ 见 \citet{attouch2016rate}。经典 Nesterov/FISTA 参数恰好是 $\alpha=3$ \citep{nesterov1983,beck2009fista}。

\paragraph{Ravine.}
外推点 $y_k$ 满足 Ravine 递推，$\gamma_k=\alpha_{k+1}$（\citealp[Theorem~2.1]{attouch2026ravine}；Ravine 方法源自 \citealp{gelfand1961}，与 Nesterov 方法的对应见 \citealp{attouch2022ravine}）。函数值与迭代点结论可以转移到 $y_k$，梯度结论则需在 $y_k$ 上单独证明。

\subsection{点一：一条能量不等式决定全部速率}

思路和 SGD 一样：找一个量，每步大致不增，然后从它的有界性读出速率。取常数步长 $s$ 来突出机制，定义能量
\[
V_k=s\,t_k^2\bigl(f(x_k)-f^\star\bigr)+\frac12\dist(z_k,S)^2 .
\]

\textbf{含义：}第一项是加权的函数值误差，权重 $t_k^2\sim k^2$ 就是加速的来源；第二项是驱动序列 $z_k$ 到解集的距离。只要 $V_k$ 有界，就有
\[
f(x_k)-f^\star\le\frac{V_k}{s\,t_k^2},
\]
速率直接由权重给出。接下来看 $V_k$ 一步之后怎么变。

利用 $z_{k+1}=z_k-s\,t_{k+1}\bigl(\nabla f(y_k)+e_k\bigr)$ 展开距离平方，再用凸性与光滑性，得到
\[
\E[V_{k+1}\mid\F_k]
\le V_k
-\underbrace{s\bigl(t_{k+1}-t_{k+1}^2+t_k^2\bigr)\bigl(f(x_k)-f^\star\bigr)}_{\text{动量条件控制}}
+\underbrace{C\,s^2t_{k+1}^2\sigma_k^2}_{\text{噪声代价}} .
\]

\textbf{含义：}能量的变化分成三部分。
\begin{itemize}
  \item 中间项比较两件事：权重从 $t_k^2$ 增长到 $t_{k+1}^2$，以及梯度步带来的下降 $t_{k+1}$。(K1) 保证下降不少于涨幅，所以这一项非正。
  \item 噪声项的来源与 SGD 附录 (A.1)--(A.2) 相同：交叉项 $\langle e_k,\cdot\rangle$ 在条件期望下为零，只剩平方项 $s^2t_{k+1}^2\|e_k\|^2$。\textbf{噪声只从这里进入。}
\end{itemize}

这个不等式正好符合 Robbins--Siegmund 引理 \citep{robbins1971} 的形式：若
\[
\E[a_{k+1}\mid\F_k]\le a_k-b_k+c_k,
\]
其中 $a_k,b_k,c_k\ge0$ 且 $\sum_k c_k<\infty$，则 $a_k$ a.s.\ 收敛，且 $\sum_k b_k<\infty$ a.s.。取 $a_k=V_k$、$c_k$ 为噪声项，得到：

\begin{itemize}
  \item \textbf{$O$ 速率.} 若
  \[
  \sum_k s_k^2t_k^2\sigma_k^2<\infty \qquad \text{(K2)},
  \]
  则 $V_k$ a.s.\ 收敛，从而有界，于是
  \[
  f(x_k)-f^\star=O\!\left(\frac{1}{s_kt_k^2}\right)\quad\text{a.s.}
  \]
  \item \textbf{$o$ 速率.} 若再有 (K1$^+$)，中间项至少是
  \[
  -(1-m)s\,t_{k+1}\bigl(f(x_k)-f^\star\bigr),
  \]
  这就是 $b_k$。于是
  \[
  \sum_k s_kt_{k+1}\bigl(f(x_k)-f^\star\bigr)<\infty\quad\text{a.s.}
  \]
  \textbf{含义（直观）：}若 $t_k^2(f(x_k)-f^\star)$ 趋于某个正数 $c$，则上面的和约为 $c\sum_k1/t_k$；当 $t_k\sim k$ 时该和发散，矛盾。所以 $O$ 可以改进为 $o$。严格证明还需要论文中的附加条件
  \[
  \sum_k\frac{t_{k+1}}{t_k^2}=\infty ,
  \]
  对 $t_k\sim k$ 自动成立。
\end{itemize}

\textbf{这一点的结论：}噪声递减时加速可以保住；持续噪声下随机 Nesterov 与 SGD 同阶。能保住多少加速，取决于 $s_kt_k\sigma_k$ 的可和性。

\paragraph{与 SGD 的对照.}
对一步不等式取期望并从 $1$ 到 $k$ 求和，中间项非正可以丢掉，得到
\[
s_kt_k^2\,\E[f(x_k)-f^\star]\le\E V_k\le V_1+C\sum_{i\le k}s_i^2t_i^2\sigma_i^2 .
\]
两边除以 $s_kt_k^2$：
\[
\E[f(x_k)-f^\star]
\lesssim
\underbrace{\frac{D^2}{s_kt_k^2}}_{\text{偏差}}
+\underbrace{\frac{\sum_{i\le k}s_i^2t_i^2\sigma_i^2}{s_kt_k^2}}_{\text{方差}} .
\]

\textbf{含义：}这与 SGD 的式 (9) 同构，只是分母从 $\sum s_t$ 变成 $s_kt_k^2$。加速放大了分母，也用同样的 $t_i^2$ 放大了噪声。

代入 $t_k\sim k$，看两种噪声体制（\citealp[Remark~3.2]{attouch2026ravine}）：

\begin{center}
\begin{tabularx}{\linewidth}{@{}l>{\raggedright\arraybackslash}Xlll@{}}
\toprule
噪声体制 & 参数 & 偏差 & 方差 & 结果 \\
\midrule
噪声递减 & $s$ 常数，$\sum_k k^2\sigma_k^2<\infty$ & $1/k^2$ & $1/k^2$ & 加速保住；加 (K1$^+$) 得 a.s.\ $o(1/k^2)$ \\
持续噪声 $\sigma_k\equiv\sigma$ & 必须 $s_k=c\,k^{-\delta}$ & $k^{\delta-2}$ & $k^{1-\delta}$（$\delta<3/2$） & 平衡于 $\delta=3/2$，期望 $O(\log k/\sqrt k)$ \\
\bottomrule
\end{tabularx}
\end{center}

持续噪声下，a.s.\ 分析还要求噪声项可和，即
\[
\sum_k k^{2-2\delta}<\infty,
\]
也就是 $\delta>3/2$。此时速率为 $k^{-(2-\delta)}$，能得到的 a.s.\ 速率都\textbf{比 $k^{-1/2}$ 略慢}，只能任意接近 $k^{-1/2}$，与 SGD 的噪声主导阶一致。

% \subsection{点二：迭代点 a.s.\ 收敛（Theorem 3.1(iii)）}

% 点一只回答``函数值多快趋于最优''，没有回答``点列是否停在某个极小点''。(i)--(ii) 给出的是
% \[
% f(x_k)\to f^\star,
% \qquad
% \|x_k-x_{k-1}\|\to0 .
% \]

% \textbf{为什么不够：}当 $S$ 不是单点时，点列可以沿着 $S$ 附近以越来越小的步子来回漂移。例如每步走 $1/k$，相邻差趋于零，但总路程发散，序列仍然可能不停下来。

% 排除这种漂移的标准工具是 Opial 引理 \citep{opial1967}。它要求两件事：
% \begin{enumerate}
%   \item 对每个 $x\in S$，$\lim_k\|x_k-x\|$ 存在；
%   \item 每个聚点都属于 $S$。
% \end{enumerate}

% \textbf{含义：}第 1 条说点列到每个解的距离最终稳定，因此不能到处漂移；第 2 条说它只能停在解上。两者合起来，聚点只能有一个。

% 四个条件分别服务于这两条：

% \begin{center}
% \begin{tabularx}{\linewidth}{@{}lcX@{}}
% \toprule
% 条件 & 服务于 & 作用 \\
% \midrule
% $\inf_k s_k>0$ & 第 2 条 & 把 $\sum s_kt_{k+1}(f-f^\star)<\infty$ 变成 $f(x_k)\to f^\star$，从而聚点是解 \\
% (K1$^+$) & 两条 & 提供可和性和 $\sum_k t_k\|x_k-x_{k-1}\|^2<\infty$ \\
% $\alpha_k\in[0,1]$ & 第 1 条 & 建立距离平方 $h_k=\tfrac12\|x_k-x\|^2$ 的准 Fejér 不等式 \\
% (K2$^+$)：$\sum_k s_kt_k\sigma_k<\infty$ & 第 1 条 & 控制该不等式中的噪声 \\
% \bottomrule
% \end{tabularx}
% \end{center}

% \textbf{为什么这里要 $\ell^1$，而点一只要 $\ell^2$：}带动量时，$h_k$ 不单调，动量可能让距离暂时变大。标准做法是证明增量的正部
% \[
% \sum_k(h_{k+1}-h_k)_+<\infty .
% \]
% 先取正部再取期望，零均值的交叉项就不能再抵消，只能用 Cauchy--Schwarz 界成一次方的 $s_kt_k\sigma_k$。所以 $\ell^1$ 来自证明手法，是否必要仍是开放问题。

% \textbf{量词问题：}对每个 $x\in S$，第 1 条成立在一个概率为 1 的事件 $\Omega_x$ 上。但 $S$ 通常不可数，不可数个概率 1 事件的交不一定概率为 1。本文取 $S$ 的可数稠密子集，在其上取交，再用逼近推广到整个 $S$，修复了 \citet{sebbouh2021} 与 \citet{liu2022} 中的这一处；同类的处理见随机准 Fejér 框架 \citep{combettes2015}。

% \textbf{与前人的关系：}

% \begin{center}
% \begin{tabularx}{\linewidth}{@{}l>{\raggedright\arraybackslash}X>{\raggedright\arraybackslash}X>{\raggedright\arraybackslash}X@{}}
% \toprule
% 工作 & 动量 & 步长 / 噪声 & 迭代点结论 \\
% \midrule
% \citet{sebbouh2021} & SHB，动量由步长决定 & $\eta_k\sim k^{-1/2-\delta}$，持续噪声 & 有限维 a.s.\ 收敛，速率与 SGD 同级 \\
% \citet{liu2022} & SHB、SNAG，常数动量 & $s_k\sim k^{-2/3-\delta}$，持续噪声 & 有限维 a.s.\ 收敛 \\
% \citet{attouch2026ravine} & 一般 $\alpha_k$，含 $\alpha_k\to1$ & $\inf s_k>0$，噪声递减 & 保留加速速率的 a.s.\ 收敛；无限维为弱收敛 \\
% \bottomrule
% \end{tabularx}
% \end{center}

% 确定性情形下，$\alpha>3$ 的迭代点收敛由 \citet{chambolle2015} 证明。在有限维中，本文的结论就是存在随机极小点 $X(\omega)\in S$，使
% \[
% \|x_k-X(\omega)\|\to0\quad\text{a.s.}
% \]

% \textbf{这一点的结论：}本文的新意在于噪声递减、保住加速的体制下得到迭代点收敛。持续噪声下，$\inf s_k>0$ 与 $\sum_k s_kt_k\sigma_k<\infty$ 不能同时成立，这一格仍然空着。

% \subsection{其余结果}

% Theorem 3.2 把函数值速率和迭代点收敛转移到 Ravine 序列 $y_k$。Theorem 3.3 在常数步长下给出
% \[
% \sum_k t_{k+1}^2\|\nabla f(y_k)\|^2<\infty\quad\text{a.s.},
% \]
% 这是 Ravine 侧独有的梯度结果，需要直接在 $y_k$ 上证明。Corollary 4.1--4.3 把一般结论代入 $1-\alpha/k$、$1-\alpha/k^r$ 和常数动量三族参数。Appendix B 用含 Hessian 阻尼项的高分辨率 ODE 解释梯度的快速收敛 \citep{shi2022highres,attouch2022ravine}。


空白的部分：

% 正文
\begin{table}[htbp]
    \centering
    \small
    \renewcommand{\arraystretch}{1.35}
    \caption{不同外推参数与噪声体制下的收敛结果}
    \label{tab:alpha-noise-regimes}

    \begin{tabularx}{\textwidth}{
        @{}
        >{\centering\arraybackslash}p{1.2cm}
        >{\raggedright\arraybackslash}X
        >{\raggedright\arraybackslash}X
        >{\raggedright\arraybackslash}X
        >{\raggedright\arraybackslash}X
        @{}
    }
        \toprule
        参数
        & 确定性
        & 随机：噪声递减
        & 随机：持续噪声
        & 随机：相对噪声 \\
        \midrule

        $\alpha=3$
        &
        函数值速率
        $\mathcal{O}(1/k^2)$；
        迭代点收敛已解决
        &
        函数值速率
        $\mathcal{O}(1/k^2)$；
        $o(1/k^2)$ 速率仍待解决；
        迭代点收敛由 Salzo (2025) 推进
        &
        期望速率接近
        $k^{-1/2}$；
        迭代点收敛尚未解决
        &
        期望速率
        $\mathcal{O}(1/k^2)$；
        a.s. 速率和迭代点收敛尚未解决
        \\

        $\alpha>3$
        &
        函数值速率
        $o(1/k^2)$；
        迭代点收敛已解决
        &
        $o(1/k^2)$、
        梯度可和性和迭代点收敛
        已由 JOTA 给出
        &
        速率接近
        $o(k^{-1/2})$；
        迭代点收敛尚未解决
        &
        JOTA 尚未覆盖
        \\

        \bottomrule
    \end{tabularx}

    \vspace{2mm}

    \begin{minipage}{0.96\textwidth}
        \footnotesize
        \textbf{说明：}
        “噪声递减”表示 $\sigma_k\to0$；
        “持续噪声”表示 $\sigma_k\asymp\sigma>0$；
        “相对噪声”表示
        $\mathbb{E}[\|e_k\|^2\mid\mathcal{F}_k]
        \le \nu^2\|\nabla f(x_k)\|^2$。
        $\alpha<3$ 不列入主表，因为它不满足 JOTA 的 $(K_1)$ 条件；
        “条件是否紧”作为备用研究问题保留。
    \end{minipage}
\end{table}

\bibliographystyle{plainnat}
\bibliography{refs}

\end{document}
